\documentclass[11pt,a4paper,fontset=none]{ctexart}

\usepackage[margin=2.25cm]{geometry}
\usepackage{amsmath,amssymb,bm}
\usepackage{booktabs}
\usepackage{longtable}
\usepackage{array}
\usepackage{xcolor}
\usepackage{hyperref}
\usepackage{enumitem}
\usepackage{caption}

\setCJKmainfont{Songti SC}
\setCJKsansfont{Songti SC}
\setCJKmonofont{Songti SC}

\hypersetup{
  colorlinks=true,
  linkcolor=blue!50!black,
  citecolor=blue!50!black,
  urlcolor=blue!50!black
}

\setlist[itemize]{leftmargin=1.7em}
\setlist[enumerate]{leftmargin=1.9em}

\title{\bfseries 超表面深度成像项目十个关键问题梳理\\
\large 从 RGB-D 仿真、PSF 计算、Fisher 信息到器件与实验路线}
\author{项目讨论整理}
\date{2026年6月}

\begin{document}
\maketitle
\tableofcontents
\newpage

\section{先把整个问题重新放到一条链路里}

目前最容易混乱的原因，是这个项目同时跨了四层东西：
\begin{enumerate}
  \item 真实世界光场如何进入光学系统；
  \item 超表面如何把不同深度、角度、波长、偏振的入射光编码成传感器图像；
  \item 如何用 Fisher information、CRLB 或互信息评价这个编码是否对深度有用；
  \item 如何用 RGB-D 数据集、深度网络和真实标定把理论链路落地。
\end{enumerate}

这四层不能混成一句 ``RGB-D 经过 PSF 卷积生成编码图''。更准确的说法应该是：
\begin{quote}
真实自然场景经过超表面相机形成 CMOS 图像，是一个高维、宽谱、空间变化、非相干辐射传输问题。开源 RGB-D 数据集没有完整光场、反射谱、偏振和相干信息，因此仿真阶段通常采用非相干、分层、空间变化 PSF 的一阶近似。这个近似不是严格真实物理，但它是 Deep Optics、Depth from Defocus with Learned Optics、PhaseCam3D、极端色散金属透镜深度成像等工作常用的可执行路线。
\end{quote}

因此，第一版项目不应该声称 ``RGB-D 完整模拟真实光场''，而应声称：
\begin{quote}
RGB-D layered PSF rendering 用于比较不同光学编码设计的相对优劣；最终可信度需要依赖真实 PSF/MTF/forward operator 标定与少量实拍校正。
\end{quote}

下面按照你提出的 10 个问题逐条回答。

\section{问题一：RGB 图和深度图如何近似进入超表面的光场}

\subsection{真实物理过程是什么}

真实场景中，每个三维点 $\bm P$ 会向相机方向发出或反射一部分光。到达超表面入瞳面的光可以写成空间、方向、波长、偏振相关的辐亮度：
\begin{equation}
  L(\bm r, \bm \omega, \lambda, p),
\end{equation}
其中 $\bm r$ 是空间位置，$\bm \omega$ 是方向，$\lambda$ 是波长，$p$ 是偏振态。理想情况下，光学系统把这个高维光场映射到传感器强度：
\begin{equation}
  I_{\mathrm{sensor}}(u,v)
  =
  \sum_p
  \int_\lambda
  \int_{\Omega}
  \mathcal T(u,v;\bm r,\bm \omega,\lambda,p)
  L(\bm r,\bm \omega,\lambda,p)
  d\bm\omega d\lambda
  + n(u,v),
\end{equation}
其中 $\mathcal T$ 是由镜头、超表面、滤光片、传感器共同决定的传输算子，$n$ 是噪声。

这个式子才是更完整的物理图像。但普通 RGB-D 数据集只给我们：
\begin{equation}
  I_{\mathrm{rgb}}(x,y),\quad D(x,y),
\end{equation}
也就是一个相机视角下的 RGB 强度图和每个像素的深度。它没有告诉我们完整光场，也没有告诉我们每个表面的真实光谱反射率、偏振态、BRDF、多路径、透明折射和镜面高光。所以我们必须近似。

\subsection{常用近似：非相干分层 PSF 渲染}

自然光下，普通漫反射场景通常可近似为非相干成像。非相干成像的核心性质是：传感器强度近似等于场景强度和强度 PSF 的卷积，而不是复振幅相干叠加。因此，如果场景中不同深度层之间相互遮挡已经由原 RGB-D 图给出，我们可以把图像按照深度切成 $K$ 个软层：
\begin{equation}
  M_k(x,y)
  =
  w_k(D(x,y)),
  \quad k=1,\ldots,K,
\end{equation}
其中 $w_k$ 是围绕深度 $z_k$ 的软分配函数，例如高斯或 triangular bin。

每一层的颜色强度为：
\begin{equation}
  I_k^c(x,y)=M_k(x,y) I_{\mathrm{rgb}}^c(x,y),
\end{equation}
其中 $c\in\{R,G,B\}$。若先忽略空间变化 PSF，则编码图可以近似为：
\begin{equation}
  I_{\mathrm{meta}}^c
  =
  \sum_{k=1}^{K}
  h^c(z_k;\phi)
  *
  I_k^c
  +
  n^c,
  \label{eq:layered_simple}
\end{equation}
其中 $h^c(z_k;\phi)$ 是给定深度 $z_k$、颜色通道 $c$ 和超表面设计 $\phi$ 时的强度 PSF。

若考虑大视场，PSF 还依赖像素位置或主光线角度 $\alpha(x,y)$：
\begin{equation}
  I_{\mathrm{meta}}^c(u,v)
  =
  \sum_{k=1}^{K}
  \sum_{x,y}
  h^c(u,v;x,y,z_k,\alpha(x,y);\phi)
  M_k(x,y) I_{\mathrm{rgb}}^c(x,y)
  +
  n^c(u,v).
  \label{eq:space_variant}
\end{equation}
这就是空间变化成像算子，不再是一个简单全局卷积。

\subsection{加入 CMOS 传感器}

进入 CMOS 之前还要考虑 sensor sampling、颜色滤波阵列、曝光、量化和噪声。更完整的离散形式是：
\begin{equation}
  \bm y
  =
  Q\left[
  \mathcal B
  \left(
  t_{\mathrm{exp}}
  \int_\lambda
  S(\lambda)
  \mathcal H_\phi(\lambda)
  [\bm X(\lambda),\bm D]
  d\lambda
  \right)
  +\bm n_{\mathrm{shot}}
  +\bm n_{\mathrm{read}}
  \right],
  \label{eq:sensor_model}
\end{equation}
其中：
\begin{itemize}
  \item $\mathcal H_\phi(\lambda)$ 是由超表面决定的空间变化成像算子；
  \item $S(\lambda)$ 包含滤光片、透镜透过率、CMOS 量子效率；
  \item $\mathcal B$ 表示 Bayer CFA 或 RGB 通道采样；
  \item $Q$ 表示模数转换和量化；
  \item $n_{\mathrm{shot}}$ 是泊松光子噪声，$n_{\mathrm{read}}$ 是读出噪声。
\end{itemize}

实际第一版可以简化为：
\begin{equation}
  y_i \sim \mathrm{Poisson}(\mu_i)+\mathcal N(0,\sigma_r^2),
\end{equation}
或者在光强足够时用高斯近似：
\begin{equation}
  \bm y=\bm \mu+\bm n,\quad \bm n\sim\mathcal N(0,\bm\Sigma).
\end{equation}

\subsection{这个近似准吗}

它不是完整真实物理，但在以下条件下是合理的一阶近似：
\begin{itemize}
  \item 自然光非相干；
  \item 大部分表面近似漫反射；
  \item 不强调透明体、镜面强反射、多路径和体散射；
  \item 目的是比较不同编码器的相对优劣；
  \item 最终会用真实 PSF 标定或真实数据校正。
\end{itemize}

它不适合严格模拟以下现象：
\begin{itemize}
  \item 透明玻璃、强镜面、金属高光；
  \item 偏振敏感材质；
  \item 宽谱反射率未知导致的颜色/深度混淆；
  \item 大景深下遮挡边界的真实散焦前后遮挡关系；
  \item 完整光场中的视差和多视角效应。
\end{itemize}

\subsection{参考文献中是否也是这样做}

是的，很多计算光学深度工作都采用类似思路：
\begin{itemize}
  \item Deep Optics / Depth from Defocus with Learned Optics 使用 RGB-D 或双像素数据，通过 depth-dependent PSF 渲染 coded image，再训练网络恢复深度；
  \item PhaseCam3D 使用 RGB-D/SceneFlow 类数据和学习相位 mask，通过深度相关 PSF 生成编码图，并讨论 CRLB/Fisher 相关指标；
  \item 3D Imaging Using Extreme Dispersion in Optical Metasurfaces 使用 SceneFlow 的 RGB 和 disparity/depth，通过分层卷积生成金属透镜编码图；
  \item 宽视场 metalens 计算成像工作常使用 spatially varying PSF 或 eigenPSF 表示，再用 deconvolution 或神经网络恢复图像。
\end{itemize}

所以这条路是有前例的，但必须在论文中清楚写出近似边界，不能把它包装成完整真实光场仿真。

\section{问题二：不同深度和角度的 PSF 如何理论计算}

\subsection{从点源到入瞳面波前}

考虑一个三维点源，位于相机坐标系中：
\begin{equation}
  \bm P=(X_0,Y_0,Z_0).
\end{equation}
超表面所在平面取 $z=0$，其孔径坐标为 $(x,y)$。点源到孔径点的距离为：
\begin{equation}
  R(x,y)
  =
  \sqrt{(x-X_0)^2+(y-Y_0)^2+Z_0^2}.
\end{equation}
若波长为 $\lambda$，波数 $k=2\pi/\lambda$，入射到超表面前的复振幅可写为：
\begin{equation}
  U_{\mathrm{in}}(x,y)
  =
  A(x,y)
  \frac{\exp\left(ikR(x,y)\right)}{R(x,y)}.
\end{equation}

在远场或小孔径近似下，点源方向由视场角 $(\alpha_x,\alpha_y)$ 表示，深度引起的曲率可以写成：
\begin{equation}
  U_{\mathrm{in}}(x,y;z,\alpha_x,\alpha_y,\lambda)
  \approx
  A(x,y)
  \exp\left[
  ik
  \left(
  x\sin\alpha_x
  +y\sin\alpha_y
  +\frac{x^2+y^2}{2z}
  \right)
  \right].
  \label{eq:incident_wave}
\end{equation}
其中线性项对应入射角，二次项对应有限距离波前曲率。长距离时，二次项非常弱，这正是远距离深度难测的物理原因。

\subsection{超表面透射函数}

若先用标量相位模型，超表面透射函数为：
\begin{equation}
  t_\phi(x,y,\lambda)
  =
  A_\phi(x,y,\lambda)
  \exp\left[i\phi(x,y,\lambda)\right].
\end{equation}
经过超表面后的场为：
\begin{equation}
  U_0(x,y)
  =
  P(x,y)
  t_\phi(x,y,\lambda)
  U_{\mathrm{in}}(x,y),
  \label{eq:after_meta}
\end{equation}
其中 $P(x,y)$ 是孔径函数。

如果要体现真正超表面，$t_\phi$ 不应只是一个相位函数，而应来自 meta-atom 响应：
\begin{equation}
  t(x,y,\lambda,\alpha,p)
  =
  A(g(x,y),\lambda,\alpha,p)
  \exp[i\phi(g(x,y),\lambda,\alpha,p)],
\end{equation}
其中 $g(x,y)$ 表示该位置 meta-atom 的几何参数，例如宽度、长度、高度、旋转角。

\subsection{传播到 CMOS：Fresnel 传播}

设超表面到传感器距离为 $d$。用 Fresnel 近似，传感器平面复振幅为：
\begin{equation}
  U_d(u,v)
  =
  \frac{e^{ikd}}{i\lambda d}
  \iint
  U_0(x,y)
  \exp\left[
  i\frac{k}{2d}
  \left((u-x)^2+(v-y)^2\right)
  \right]
  dxdy.
  \label{eq:fresnel}
\end{equation}
对应强度 PSF 为：
\begin{equation}
  h(u,v;z,\alpha_x,\alpha_y,\lambda,\phi)
  =
  |U_d(u,v)|^2.
\end{equation}
通常还要归一化：
\begin{equation}
  \sum_{u,v}h(u,v;z,\alpha,\lambda,\phi)=\eta(z,\alpha,\lambda),
\end{equation}
其中 $\eta$ 表示总透过效率。

\subsection{Fraunhofer / Fourier 形式}

如果传感器位于透镜焦平面附近，或者系统可等效为 pupil-to-focal-plane Fourier transform，则 PSF 可以写成：
\begin{equation}
  h(u,v)
  =
  \left|
  \mathcal F
  \left\{
  P(x,y)
  \exp[i\Phi_{\mathrm{total}}(x,y)]
  \right\}
  \right|^2,
  \label{eq:fourier_psf}
\end{equation}
其中总相位包含：
\begin{equation}
  \Phi_{\mathrm{total}}
  =
  \phi_{\mathrm{meta}}(x,y,\lambda)
  +
  \phi_{\mathrm{lens}}(x,y,\lambda)
  +
  k(x\sin\alpha_x+y\sin\alpha_y)
  +
  k\frac{x^2+y^2}{2z}
  +
  \phi_{\mathrm{aberr}}(x,y,\alpha,\lambda).
\end{equation}
这个形式适合快速仿真和优化。

\subsection{矢量超表面形式}

真正的超表面可能对偏振有耦合。此时透射不再是标量，而是 Jones 矩阵：
\begin{equation}
  \bm J(x,y,\lambda,\alpha)
  =
  \begin{bmatrix}
  t_{xx} & t_{xy}\\
  t_{yx} & t_{yy}
  \end{bmatrix}.
\end{equation}
输入电场为 $\bm E_{\mathrm{in}}=[E_x,E_y]^T$，输出为：
\begin{equation}
  \bm E_0(x,y)=\bm J(x,y,\lambda,\alpha)\bm E_{\mathrm{in}}(x,y).
\end{equation}
然后 $E_x,E_y$ 分别传播到传感器，强度为：
\begin{equation}
  h(u,v)=|E_x(u,v)|^2+|E_y(u,v)|^2,
\end{equation}
若传感器或偏振片选择某一偏振，则还要乘相应分析矩阵。

\subsection{实际计算时能否用 RCWA}

可以，但要分层使用。

RCWA 适合计算单个周期性 meta-atom 在不同波长、入射角、偏振下的复透射：
\begin{equation}
  g=(w,l,h,\theta)
  \mapsto
  \bm J_g(\lambda,\alpha).
\end{equation}
然后建立一个 meta-atom library：
\begin{equation}
  \mathcal L=\{g_m,\ A_m(\lambda,\alpha,p),\ \phi_m(\lambda,\alpha,p)\}_{m=1}^{M}.
\end{equation}

但 RCWA 直接仿真整个毫米级甚至厘米级超表面几乎不可行，因为周期数太多。实际常用路线是：
\begin{enumerate}
  \item 用 RCWA/FDTD 仿真局部周期单元，得到 meta-atom library；
  \item 用局部周期近似 LPA，把每个位置映射到一个 meta-atom；
  \item 得到整个孔径的复透射函数 $t(x,y,\lambda,\alpha,p)$；
  \item 用 Fourier optics / Fresnel propagation 计算系统 PSF；
  \item 对关键区域或小尺寸样片用全波仿真校验 LPA 误差。
\end{enumerate}

所以回答是：RCWA 很重要，但它主要用于 meta-atom 层；系统级 PSF 计算通常靠 Fourier optics、角谱法或光线-波动混合模型。

\section{问题三：Fisher 是什么，为什么不能只对 PSF 做 Fisher}

\subsection{Fisher information 的定义}

如果观测 $\bm y$ 的概率分布为 $p(\bm y|\theta)$，参数为 $\theta$，Fisher information 定义为：
\begin{equation}
  \mathcal I(\theta)
  =
  \mathbb E_{\bm y|\theta}
  \left[
  \left(
  \frac{\partial}{\partial\theta}
  \log p(\bm y|\theta)
  \right)^2
  \right].
\end{equation}
多参数情形为：
\begin{equation}
  F_{ij}
  =
  \mathbb E
  \left[
  \frac{\partial \log p(\bm y|\bm\theta)}{\partial \theta_i}
  \frac{\partial \log p(\bm y|\bm\theta)}{\partial \theta_j}
  \right].
\end{equation}

如果噪声是高斯：
\begin{equation}
  \bm y\sim\mathcal N(\bm\mu(\bm\theta),\bm\Sigma),
\end{equation}
且协方差不依赖参数，则：
\begin{equation}
  \bm F
  =
  \bm J^T\bm\Sigma^{-1}\bm J,
  \quad
  \bm J=\frac{\partial \bm\mu}{\partial \bm\theta}.
\end{equation}

如果是泊松噪声：
\begin{equation}
  y_i\sim\mathrm{Poisson}(\mu_i(\bm\theta)),
\end{equation}
则：
\begin{equation}
  F_{ij}
  =
  \sum_m
  \frac{1}{\mu_m}
  \frac{\partial\mu_m}{\partial\theta_i}
  \frac{\partial\mu_m}{\partial\theta_j}.
\end{equation}

\subsection{PSF 级 Fisher}

PSF 级 Fisher 把一个点源的传感器响应当成观测：
\begin{equation}
  \bm y=\bm h(z,\alpha,\lambda;\phi)+\bm n.
\end{equation}
如果只估计深度：
\begin{equation}
  \mathcal I_z^{\mathrm{psf}}
  =
  \left(
  \frac{\partial \bm h}{\partial z}
  \right)^T
  \bm\Sigma^{-1}
  \left(
  \frac{\partial \bm h}{\partial z}
  \right).
\end{equation}
这衡量 ``单个点目标的 PSF 随深度变化是否明显''。

如果同时考虑深度、角度、亮度、波长等参数：
\begin{equation}
  \bm\theta=[z,\alpha_x,\alpha_y,a,\lambda]^T.
\end{equation}
此时真正关心的不是 $F_{zz}$，而是扣除 nuisance 混淆后的有效深度信息：
\begin{equation}
  \mathcal I_z^{\mathrm{eff}}
  =
  F_{zz}-\bm F_{z\eta}\bm F_{\eta\eta}^{-1}\bm F_{\eta z}.
\end{equation}

\subsection{为什么 PSF 级 Fisher 不够}

你这个怀疑完全正确。PSF 随深度变化明显，不等于整张自然图像中的深度可恢复。原因包括：
\begin{itemize}
  \item 自然图像不是孤立点源，而是纹理、边界、遮挡和多个深度层的叠加；
  \item 某些深度变化只在高频纹理上可见，低纹理区域仍然不可分；
  \item PSF 变化可能主要表现为模糊，但模糊也可能来自运动、曝光、材质和噪声；
  \item 宽视场下角度变化可能和深度变化混淆；
  \item 多个深度层卷积后可能互相抵消或遮挡。
\end{itemize}

因此应该分三级：
\begin{enumerate}
  \item PSF 级 Fisher：验证物理编码器有没有局部深度敏感性；
  \item 层状场景 Fisher：对多个深度层、局部纹理 patch 计算观测图像的 Fisher；
  \item 图像级/任务级 Fisher 或互信息：直接衡量整张编码图对深度图的有效信息。
\end{enumerate}

\subsection{图像级 Fisher 如何写}

设 RGB-D 场景为 $\bm X,\bm D$，经过编码器后传感器均值为：
\begin{equation}
  \bm\mu_\phi=\mathcal H_\phi(\bm X,\bm D).
\end{equation}
如果把某个像素或某个深度参数 $D_j$ 作为待估参数，高斯噪声下：
\begin{equation}
  F_{ij}^{\mathrm{img}}
  =
  \left(
  \frac{\partial \bm\mu_\phi}{\partial D_i}
  \right)^T
  \bm\Sigma^{-1}
  \left(
  \frac{\partial \bm\mu_\phi}{\partial D_j}
  \right).
\end{equation}
完整 Fisher matrix 的维度等于像素数，极其巨大。实际不能直接完整求逆。可行近似包括：
\begin{itemize}
  \item patch-level Fisher：每次只对局部 patch 的几个深度参数计算；
  \item random projection / Hutchinson trace estimator：估计 $\mathrm{Tr}(F)$ 或 $\mathrm{Tr}(F^{-1})$；
  \item 只对低维深度参数计算，例如平面深度、局部尺度、远距离 bin；
  \item 用互信息下界或对比学习替代直接 Fisher；
  \item 直接用端到端 depth loss 作为任务级指标。
\end{itemize}

\subsection{相关工作是否已有}

已有工作确实使用 Fisher/CRLB 分析或优化相位 mask、PSF 和深度估计。PhaseCam3D 是最相关之一，它使用 learned phase masks、PSF 仿真和深度网络，并讨论 Fisher/CRLB 对深度相关 PSF 的指导作用。Depth from Defocus with Learned Optics、Deep Optics 系列也有类似 ``学习光学编码 + 深度网络'' 的思想。

但你的潜在区别不应是 ``第一次把 Fisher 用到相位 mask''，而应是：
\begin{quote}
把 Fisher/CRLB 从中心视场或简单深度 PSF 扩展到长距离、大视场、宽谱、角度/光谱/亮度 nuisance-aware 的有效深度信息，并进一步加入图像级 Fisher/互信息和真实场景任务损失。
\end{quote}

所以不要把创新写成 ``提出 Fisher 优化 PSF''。那太危险。更稳的写法是：
\begin{quote}
提出面向长距离宽视场自然场景的多参数有效 Fisher 深度编码设计，并系统研究 PSF 级信息、图像级信息和最终深度任务表现之间的一致性与差异。
\end{quote}

\section{问题四：自然光自然场景下波长怎么处理}

\subsection{最完整形式}

自然光图像应该对波长积分：
\begin{equation}
  I_{\mathrm{meta}}(u,v)
  =
  \int_{\lambda_{\min}}^{\lambda_{\max}}
  S_{\mathrm{illum}}(\lambda)
  R_{\mathrm{scene}}(x,y,\lambda)
  T_{\mathrm{optics}}(\lambda)
  H_\phi(u,v;x,y,D,\lambda)
  d\lambda.
\end{equation}
其中 $S_{\mathrm{illum}}$ 是照明光谱，$R_{\mathrm{scene}}$ 是物体反射谱，$T_{\mathrm{optics}}$ 是系统透过率。

普通 RGB 图只给三个通道：
\begin{equation}
  I^c(x,y)=
  \int S_c(\lambda)
  S_{\mathrm{illum}}(\lambda)
  R_{\mathrm{scene}}(x,y,\lambda)
  d\lambda,
\end{equation}
并不能反推出完整 $R_{\mathrm{scene}}(\lambda)$。所以自然光宽谱仿真有天然不确定性。

\subsection{三种可执行路线}

\paragraph{路线 A：窄带滤波}
在超表面前加入窄带滤光片，例如中心波长 $\lambda_0$、带宽 10--30 nm。此时近似：
\begin{equation}
  I_{\mathrm{meta}}\approx H_\phi(\lambda_0)[I,D].
\end{equation}
优点是仿真和硬件简单，PSF 清晰可控；缺点是光通量下降，长距离 SNR 变差，颜色信息损失。

\paragraph{路线 B：RGB 三波段近似}
把 RGB 通道当成三个宽波段，分别用代表波长 $\lambda_R,\lambda_G,\lambda_B$：
\begin{equation}
  I_{\mathrm{meta}}^c
  =
  \mathcal H_\phi(\lambda_c)[I^c,D].
\end{equation}
这是很多早期仿真最容易做的近似。优点是可以直接使用 RGB-D 数据；缺点是忽略通道内宽谱积分。

\paragraph{路线 C：多波长采样积分}
在每个颜色通道内采样多个波长：
\begin{equation}
  I_{\mathrm{meta}}^c
  =
  \sum_{\ell=1}^{L}
  w_{c,\ell}
  \mathcal H_\phi(\lambda_{c,\ell})[\tilde I_{c,\ell},D].
\end{equation}
由于没有真实光谱，可以先令 $\tilde I_{c,\ell}\approx I^c$ 或加入随机光谱扰动。这个方法更接近自然光，但不确定性也更大。

\subsection{建议第一版怎么做}

第一版建议同时做两个版本：
\begin{enumerate}
  \item 单波长/窄带版本：验证理论和 Fisher 优化是否成立；
  \item RGB 三波段或多波长版本：测试自然光宽谱是否把编码冲淡。
\end{enumerate}

如果单波长下有效、多波长下失效，这不是失败，而是重要发现。它会告诉你系统必须使用滤波、窄带照明、多通道复用或色散工程。

\section{问题五：整个工作本质上涉及哪些光传播数值理论}

这个项目确实很大程度是 ``光传播 + 统计估计 + 学习重建'' 的数值理论计算。核心模块如下。

\begin{longtable}{p{3.2cm}p{5.4cm}p{6.2cm}}
\toprule
模块 & 作用 & 常用方法 \\
\midrule
\endhead
点源到孔径波前 & 描述深度和视场角如何进入系统 & 球面波、平面波近似、Fresnel 近似 \\
\addlinespace
超表面局部响应 & 由 meta-atom 几何得到相位、振幅、偏振、色散 & RCWA、FDTD、FEM、局部周期近似 \\
\addlinespace
孔径到传感器传播 & 由透射场计算 PSF & Fourier optics、Fresnel propagation、angular spectrum method \\
\addlinespace
空间变化成像 & 处理大视场下不同像素不同 PSF & spatially varying convolution、eigenPSF、patch-wise convolution \\
\addlinespace
宽谱积分 & 自然光、多波段、滤光片建模 & 波长采样、RGB 通道近似、光谱随机化 \\
\addlinespace
CMOS 模型 & 把光强转为数字图像 & shot noise、read noise、Bayer CFA、gamma/linear RAW、ADC \\
\addlinespace
Fisher/CRLB & 评价理论可估计性 & Gaussian/Poisson Fisher、多参数 Schur complement \\
\addlinespace
图像级信息 & 评价自然场景中是否有用 & patch Fisher、互信息下界、InfoNCE、task loss \\
\addlinespace
学习重建 & 从编码图恢复深度 & CNN/Transformer decoder、teacher distillation、uncertainty \\
\bottomrule
\end{longtable}

因此，你可以把代码分成四个包：
\begin{enumerate}
  \item \texttt{optics}: PSF 和传播计算；
  \item \texttt{render}: RGB-D 到编码图；
  \item \texttt{info}: Fisher/CRLB/互信息指标；
  \item \texttt{learn}: depth decoder 和训练评估。
\end{enumerate}

\section{问题六：超表面如何设计，口径、周期、尺寸和可制造性}

\subsection{口径决定长距离信息上限}

长距离深度的物理难点来自波前曲率变化很小。对孔径半径 $a$ 和距离 $z$，有限距离带来的孔径边缘相位量级近似为：
\begin{equation}
  \Delta\Phi_z
  \sim
  k\frac{a^2}{2z}.
\end{equation}
两段距离 $z$ 和 $z+\Delta z$ 的相位差约为：
\begin{equation}
  \Delta\Phi
  \sim
  k\frac{a^2}{2}
  \left|
  \frac{1}{z}-\frac{1}{z+\Delta z}
  \right|
  \approx
  k\frac{a^2\Delta z}{2z^2}.
\end{equation}
这说明：
\begin{itemize}
  \item 深度敏感性随口径平方增加；
  \item 随距离平方下降；
  \item 小口径在远距离上天然非常吃亏。
\end{itemize}

所以必须先做 feasibility map：给定口径、波长、像素尺寸、曝光、FOV，理论上 20 m、50 m、100 m 能分辨多少深度差。没有这个图，不应该承诺长距离高精度。

\subsection{大视场要求角度响应可控}

大视场并不是 ``超表面做大一点''。大视场要求离轴入射时：
\begin{equation}
  t(x,y,\lambda,\alpha)
\end{equation}
仍然可控。普通局部相位设计在大角度下会出现：
\begin{itemize}
  \item 离轴像差；
  \item 焦点漂移；
  \item 效率下降；
  \item 偏振串扰；
  \item 角度和深度编码混淆。
\end{itemize}

因此超表面设计要么使用角度鲁棒 meta-atom，要么承认 PSF 是空间变化的，并把 $\alpha$ 输入到 forward model 和 decoder。

\subsection{周期和采样}

meta-atom 周期 $p$ 通常要小于波长量级，避免高阶衍射：
\begin{equation}
  p < \frac{\lambda}{n_{\mathrm{env}}(1+\sin\alpha_{\max})}
\end{equation}
这是粗略条件，具体与材料和结构有关。可见光下周期常在几百纳米量级。若口径为 $D_{\mathrm{ap}}$，单元数约为：
\begin{equation}
  N\approx \left(\frac{D_{\mathrm{ap}}}{p}\right)^2.
\end{equation}
例如 $D_{\mathrm{ap}}=1$ mm，$p=400$ nm，则 $N\approx6.25\times10^6$ 个单元；若 $D_{\mathrm{ap}}=1$ cm，则 $N\approx6.25\times10^8$，制造和设计压力巨大。

\subsection{第一版设计建议}

不要一开始直接设计厘米级全功能超表面。建议路线：
\begin{enumerate}
  \item 从等效 phase mask 做理论；
  \item 加相位量化，例如 8/16/32 levels；
  \item 加透过率约束；
  \item 用公开或自己 RCWA 生成的小型 meta-atom library 映射；
  \item 先考虑毫米级样片或混合光学结构；
  \item 若要长距离，优先考虑普通镜头 + 超表面编码片，而不是单片超表面替代所有光学功能。
\end{enumerate}

\subsection{小超表面是否天生无法接受大视场}

小超表面不是不能接受大视场，而是存在两个问题：
\begin{itemize}
  \item 小口径导致长距离深度信息和光通量不足；
  \item 大视场下入射角范围大，meta-atom 响应更容易失效。
\end{itemize}
因此小超表面可以做宽视场编码演示，但若目标是长距离高精度 metric depth，必须认真评估口径、SNR 和视场角。很可能更现实的架构是：
\begin{quote}
普通成像镜头提供口径和基本成像能力，超表面作为 pupil-plane 或 sensor-near 编码器提供额外深度信息。
\end{quote}

\section{问题七：还有哪些没有考虑到的问题}

除了你已经列出的点，还要特别注意以下问题。

\subsection{线性 RAW 与普通 RGB 图像}

数据集 RGB 往往经过 ISP、gamma、白平衡、压缩，而物理成像模型需要线性辐照度。若直接用 sRGB 做卷积，严格来说不物理。第一版可以接受，但要说明；更好是尽量使用 linear RGB 或对 sRGB 做 inverse gamma。

\subsection{景深、对焦和平面位置}

深度相关 PSF 取决于传感器位置、焦距、基准对焦距离。必须明确系统是：
\begin{itemize}
  \item focused at infinity；
  \item focused at some distance；
  \item coded aperture near pupil；
  \item metalens directly focuses to sensor。
\end{itemize}
否则同一个深度 $z$ 的 PSF 没有定义。

\subsection{遮挡边界}

layered convolution 在深度边界处会有前景/背景混合问题。真实散焦时前景会遮挡背景光，不是简单所有层线性相加。需要使用 soft matting 或 occlusion-aware rendering，否则边界可能不准。

\subsection{动态范围和曝光}

超表面编码可能把光能分散到大范围 PSF，导致峰值变暗、读噪增加、动态范围不足。Fisher 目标必须考虑 photon budget，而不是只看 PSF 变化。

\subsection{真实标定漂移}

温度、装调、传感器 cover glass、滤光片角度、封装误差都会改变 PSF。需要把 calibration code 或 PSF library 输入网络。

\subsection{评价目标}

无人机不一定需要每个像素厘米级深度。可能更现实的目标是：
\begin{itemize}
  \item 远距离障碍物距离分段；
  \item metric scale anchor；
  \item 低纹理区域相对纯 RGB 的改善；
  \item 边缘视场障碍物检测深度可靠性。
\end{itemize}

\section{问题八：虚拟数据算 Fisher 时，视野、景深、相机参数是否关键}

非常关键。RGB-D 数据不是随便拿来就能进你的超表面模型。必须把数据集相机映射到你的目标相机。

\subsection{需要的相机信息}

至少需要：
\begin{itemize}
  \item 相机内参 $f_x,f_y,c_x,c_y$；
  \item 分辨率；
  \item FOV；
  \item depth 定义，是 optical axis depth 还是 range depth；
  \item depth 单位和有效范围；
  \item RGB 是否线性；
  \item 是否有畸变；
  \item 是否有相机位姿，若做视频自监督。
\end{itemize}

每个像素对应的 ray direction 是：
\begin{equation}
  \bm r(x,y)
  =
  \frac{
  K^{-1}[x,y,1]^T
  }{
  \|K^{-1}[x,y,1]^T\|
  }.
\end{equation}
视场角 $\alpha(x,y)$ 来自这个 ray direction。

\subsection{为什么必须清洗}

如果数据集原始 FOV 是 60 度，而你仿真目标是 120 度，直接拉伸图像是错误的。因为每个像素对应的入射角不同，PSF 也不同。必须做：
\begin{itemize}
  \item 按目标相机内参重新投影；
  \item 或只使用数据集中与目标 FOV 匹配的相机；
  \item 或明确仿真的 FOV 就等于数据集 FOV。
\end{itemize}

\subsection{建议数据策略}

第一版选择有明确内参和 dense depth 的合成数据，例如 TartanAir、Virtual KITTI、Hypersim。处理流程：
\begin{enumerate}
  \item 读取 RGB、depth、intrinsics；
  \item 转成线性 RGB；
  \item 根据 intrinsics 计算每个像素 ray angle；
  \item 裁剪或重投影到目标 FOV；
  \item 按目标距离范围 mask；
  \item 再进入 PSF rendering。
\end{enumerate}

这样你的 Fisher 和训练图像才对应真实系统几何。

\section{问题九：理想实验和现实实验}

\subsection{理想实验}

理想实验确实是：
\begin{equation}
  \text{真实场景}
  \rightarrow
  \text{超表面}
  \rightarrow
  \text{CMOS}
  \rightarrow
  \text{编码图}
  \rightarrow
  \text{深度重建}.
\end{equation}
同时用 LiDAR、ToF 或高精度双目系统提供真值深度。

\subsection{现实中通常还需要什么}

现实系统很可能需要：
\begin{itemize}
  \item 普通镜头：提供口径、聚焦和成像基础；
  \item 超表面编码片：放在 pupil plane、image plane 附近或 sensor-near；
  \item 窄带滤光片或多窄带滤光片：控制波长；
  \item 偏振片：读取偏振复用通道；
  \item 标定板和点光源：测 PSF/MTF；
  \item 参考相机或分光器：采集 paired RGB；
  \item LiDAR/ToF：提供 metric depth 真值；
  \item 受控 LED/激光：在早期实验中减少自然光不确定性。
\end{itemize}

\subsection{是否用激光/LED}

如果目标是被动自然光，最终应该避免主动照明。但早期验证可以用 LED 或窄带光源，因为它能先证明光学编码是否符合设计。合理路线是：
\begin{enumerate}
  \item 窄带 LED + 点源/标定板测 PSF；
  \item 窄带 LED + 简单深度物体验证重建；
  \item 室内自然光/宽带光测试；
  \item 户外自然光测试。
\end{enumerate}

这不是退步，而是工程上降低不确定性的正常方法。

\section{问题十：建议的完整技术路线}

\subsection{第一阶段：物理可行性}

目标：先证明长距离大视场下是否有足够深度信息。

要做：
\begin{enumerate}
  \item 设定目标系统：波长、口径、FOV、像素尺寸、传感器距离、目标深度范围；
  \item 实现点源 PSF 计算：$h(z,\alpha,\lambda;\phi)$；
  \item 计算普通透镜、人工 PSF、随机相位的 CRLB；
  \item 优化低维相位，使远距离/边缘视场 $\mathcal I_z^{\mathrm{eff}}$ 提高；
  \item 输出 feasibility map。
\end{enumerate}

如果这一步显示 50--100 m 的 CRLB 天然很差，就及时调整目标，例如改成 5--30 m 或远距离尺度锚点。

\subsection{第二阶段：RGB-D 编码图仿真}

目标：验证 PSF 级优势是否转化为自然图像优势。

要做：
\begin{enumerate}
  \item 选择 TartanAir / Virtual KITTI / Hypersim；
  \item 清洗 intrinsics、FOV、depth 单位；
  \item 用 layered spatially varying PSF 渲染编码图；
  \item 比较不同相位设计下编码图的图像级 Fisher 或局部可辨性；
  \item 训练小型 decoder。
\end{enumerate}

\subsection{第三阶段：严格消融}

必须比较：
\begin{itemize}
  \item RGB-only；
  \item Meta-only；
  \item RGB+Meta；
  \item RGB+Meta+ray angle；
  \item random phase；
  \item fixed DH/rotating PSF；
  \item Fisher-only phase；
  \item Fisher + image/task phase。
\end{itemize}

分段报告：
\begin{itemize}
  \item 近/中/远距离；
  \item 中心/边缘视场；
  \item 高纹理/低纹理；
  \item 单波长/宽谱；
  \item 高 SNR/低 SNR。
\end{itemize}

\subsection{第四阶段：metasurface 约束}

在理论有效后，再引入：
\begin{itemize}
  \item phase quantization；
  \item amplitude/efficiency；
  \item wavelength dependence；
  \item polarization multiplexing；
  \item RCWA meta-atom library；
  \item fabrication perturbation。
\end{itemize}

\subsection{第五阶段：实验原型}

先从简单原型开始：
\begin{enumerate}
  \item SLM 或 DOE 验证相位编码；
  \item 测 PSF library；
  \item 用真实 PSF 替代理想 PSF 训练；
  \item 再考虑 metasurface fabrication；
  \item 若要户外长距离，优先考虑 hybrid lens + metasurface encoder。
\end{enumerate}

\section{最核心的结论}

现在这个项目可以压缩成以下判断：
\begin{enumerate}
  \item RGB-D + PSF 卷积不是完整光场仿真，但它是自然光非相干条件下的合理一阶模型，且已有计算光学深度工作大量采用；
  \item 点源 PSF 可以由球面波入射、超表面复透射和 Fresnel/Fourier 传播理论计算；
  \item RCWA 适合做 meta-atom library，不适合直接仿真整个大口径系统；
  \item Fisher 不能只看 PSF，必须从 PSF 级扩展到 nuisance-aware、patch/image-level 和 task-level；
  \item 自然光波长处理是大风险，第一版应同时做窄带理想验证和宽谱退化分析；
  \item 长距离成败首先由口径、SNR、FOV 和波长决定，必须先做 feasibility map；
  \item 数据集必须按相机内参、FOV、depth 定义清洗，否则仿真几何不成立；
  \item 真实实验大概率需要普通镜头、滤光片、偏振片、标定系统或受控光源，不能想象成裸超表面直接解决一切；
  \item 最稳路线是先做等效 phase mask 的信息最优编码，再逐步加入 metasurface 的色散、偏振、角度和制造约束。
\end{enumerate}

\section*{参考方向}

建议重点阅读或复查以下方向：
\begin{itemize}
  \item PhaseCam3D: Learning Phase Masks for Passive Single View Depth Estimation；
  \item Depth from Defocus with Learned Optics；
  \item Deep Optics for Monocular Depth Estimation and 3D Object Detection；
  \item 3D Imaging Using Extreme Dispersion in Optical Metasurfaces；
  \item Monocular metasurface camera for passive single-shot 4D imaging；
  \item wide field-of-view metalens imaging with deep learning；
  \item thin on-sensor nanophotonic array cameras；
  \item RCWA / local periodic approximation based metalens design。
\end{itemize}

\end{document}
